\chapter{Global Sensitivity Analysis and Sobol Indices}

\begin{flushright}
\textit{The point is not to remove uncertainty.\\
The point is to find out which uncertainty is worth worrying about.}
\end{flushright}

\bigskip

Consider two epidemiological models. Both give $\SI{\beta}{R_0} = +1$
at the same nominal parameter values. In Model~A, $R_0$ is nearly linear
in $\beta$ across the full biologically plausible range $[0.01, 0.15]$.
In Model~B, $R_0$ is highly nonlinear: it is steep near the lower end
of the range and flat at the upper end.

For Model~A, the local SI correctly identifies $\beta$ as an important
parameter. For Model~B, local SA gives the same SI value of $+1$;
but it is now describing only the slope at the nominal point.
A $20\%$ increase in $\beta$ from the nominal produces very different
effects in the two models, and no derivative computed at the nominal
value can see this.

Chapter~8 showed when local SA fails. This chapter provides the tools
to replace it.

%%------------------------------------------------------------
\section{Why Local SA Is Not Enough: The Sigma-Normalized Bridge}
\label{sec:sigmanorm}
%%------------------------------------------------------------

\begin{flushright}
\textit{\small Local SA studies the neighborhood. Global SA studies the country.}
\end{flushright}

The two can give very different answers about which parameters matter most,
as the following example shows.

\subsection{The Portfolio Example}

Consider a model $Y = s_1 Q_1 + s_2 Q_2$ where $s_1 = 2$, $s_2 = 1$
and the parameters have standard deviations $\sigma_1 = 1$, $\sigma_2 = 3$.

Local SA gives $\SI{Q_1}{Y} = 2 > \SI{Q_2}{Y} = 1$:
parameter~1 appears more important. But the parameters have very different
uncertainty ranges, and $Q_2$ varies three times as much as $Q_1$.

The output variance is:
\[
  \Var(Y) = s_1^2\sigma_1^2 + s_2^2\sigma_2^2 = 4 + 9 = 13
  \quad\implies\quad \sigma_Y = \sqrt{13}.
\]
The \textbf{sigma-normalized sensitivity index}\index{sigma-normalized index}
\begin{equation}
  S_i^\sigma = \frac{\sigma_i}{\sigma_Y}\,\pd{Y}{Q_i}
  \label{eq:sigmanorm}
\end{equation}
accounts for both the derivative and the parameter uncertainty:
\[
  S_1^\sigma = \frac{1}{\sqrt{13}}\cdot 2 \approx 0.55, \qquad
  S_2^\sigma = \frac{3}{\sqrt{13}}\cdot 1 \approx 0.83.
\]
Parameter~2 dominates, opposite to the local SA ranking.
The local ranking was misleading because it ignored the factor-of-3
difference in parameter uncertainty.

\unifyingidea{1} The sigma-normalized index is still the sensitivity
index from Chapter~2, multiplied by the relative parameter uncertainty.
Global SA builds on local SA; it does not replace it. The local SI
from Chapter~2 appears inside every global SA formula in this chapter.

\subsection{Two SIR Models}

For the SIR model, suppose two outbreak settings give the same
$\SI{\beta}{R_0} = 1$. In Setting~A, $\beta$ is tightly constrained
from contact-tracing data ($\sigma_\beta = 0.005$), while in Setting~B,
$\beta$ must be estimated from seroprevalence data ($\sigma_\beta = 0.03$).
Local SA says $\beta$ is equally important in both settings.
The sigma-normalized index says $\beta$ in Setting~B is six times more
important, in terms of its contribution to uncertainty in $R_0$.
The second semester will show how to propagate these uncertainties
through the full model.

\begin{selfcheck}
For the SIR model with local SIs $\SI{k}{R_0} = \SI{\beta}{R_0} = \SI{\tau}{R_0} = 1$
and parameter standard deviations $\sigma_k = 0.5$, $\sigma_\beta = 0.01$,
$\sigma_\tau = 2$, with $\sigma_{R_0} = \sqrt{\sigma_k^2 + \sigma_\beta^2 + \sigma_\tau^2}/R_0^*$:

For a SIR model with $k^* = 5$, $\beta^* = 0.06$, $\tau^* = 7$
($R_0^* = 2.1$), first compute $\sigma_{R_0}$ from the delta method, then
rank the three parameters by their sigma-normalized index $S_i^\sigma$.
\end{selfcheck}
\begin{scAnswer}
Delta method: $\Var(R_0) \approx \sum_i (\partial R_0/\partial p_i)^2\sigma_i^2$.
$\partial R_0/\partial k = \beta\tau = 0.42$; $\partial R_0/\partial\beta = k\tau = 35$;
$\partial R_0/\partial\tau = k\beta = 0.3$.
$\Var(R_0) \approx 0.42^2(0.25) + 35^2(0.0001) + 0.3^2(4)
= 0.0441 + 0.1225 + 0.36 = 0.527$.
$\sigma_{R_0} \approx 0.726$.

$S_k^\sigma = (0.5/0.726)(0.42/2.1)(k/k)\cdot$... Easier via the formula directly:
$S_i^\sigma = (\partial R_0/\partial p_i)\sigma_i/\sigma_{R_0}$.
$S_k^\sigma = 0.42 \times 0.5/0.726 \approx 0.29$.
$S_\beta^\sigma = 35 \times 0.01/0.726 \approx 0.48$.
$S_\tau^\sigma = 0.3 \times 2/0.726 \approx 0.83$.

Ranking: $\tau > \beta > k$. The mean infectious period drives the most
uncertainty in $R_0$, despite all three having equal local SIs.
\end{scAnswer}

%%------------------------------------------------------------
\section{Sampling Methods}
\label{sec:sampling}
%%------------------------------------------------------------

All global SA methods require evaluating the model at many parameter
combinations. Choosing those combinations efficiently is the first
computational task.

\subsection{Random Sampling}

The simplest approach: draw $N$~independent samples from the joint
parameter distribution (often assumed uniform or independent log-normal
over plausible ranges). For $N$ large, random sampling provides good
coverage on average, but can leave regions of parameter space unsampled
by chance.

\subsection{Latin Hypercube Sampling}

Latin Hypercube Sampling (LHS) divides each parameter's range into
$N$ equal strata and ensures that exactly one sample falls in each
stratum for every parameter~\cite{mckay1979comparison}. This eliminates the clustering that
random sampling can produce and improves coverage with smaller $N$.

\begin{definition}[LHS Sample]\index{Latin hypercube sampling (LHS)}
An LHS of size $N$ for $m$ parameters with ranges $[l_j, u_j]$,
$j = 1,\ldots,m$, is an $N\times m$ matrix $X$ such that each column
contains exactly one value from each of the $N$ equal-width strata
$[l_j + (k-1)(u_j-l_j)/N,\ l_j + k(u_j-l_j)/N]$ for $k = 1,\ldots,N$.
\end{definition}

The MATLAB implementation generates LHS by permuting the strata
independently for each parameter and adding uniform jitter within each stratum:

\begin{verbatim}
%% From: src/globalsa/lhs_sample.m (see GitHub repository)
rng(42);                                   %% set seed before sampling
X = lhs_sample(N, m, lb, ub);             %% N-by-m sample matrix
%% Uses lhsdesign (Statistics Toolbox) with maximin criterion
\end{verbatim}

For the same $N$, LHS provides substantially better coverage than
random sampling, particularly for estimating marginal effects of
individual parameters; see Helton and Davis~\cite{helton2003latin}
for convergence analysis and uncertainty propagation results.
For Sobol index estimation (Section~\ref{sec:sobol}),
LHS is not used directly; instead, the Saltelli quasi-random construction
is preferred.

\begin{selfcheck}
For $N = 4$ samples and $m = 1$ parameter on $[0, 1]$:
(a) Draw a random sample of size 4. Is it possible that all 4 points
    fall in $[0, 0.5]$? Is this possible with LHS?
(b) Write out explicitly what an LHS of size 4 looks like for one parameter.
\end{selfcheck}
\begin{scAnswer}
(a) Yes, random sampling can place all 4 points in $[0, 0.5]$ (probability $(1/2)^4 = 1/16$).
With LHS, this is impossible: each quarter-interval $[0,0.25)$, $[0.25,0.5)$,
$[0.5,0.75)$, $[0.75,1]$ must contain exactly one sample.

(b) One possible LHS: one point in each of $[0,0.25)$, $[0.25,0.5)$,
$[0.5,0.75)$, $[0.75,1]$, with uniform jitter within each. Example:
$X = \{0.18, 0.43, 0.62, 0.91\}$, one value per stratum,
randomly positioned within each.
\end{scAnswer}

%%------------------------------------------------------------
\section{Regression-Based SA: LHS/PRCC}
\label{sec:prcc}
%%------------------------------------------------------------

The Partial Rank Correlation Coefficient (PRCC)\index{PRCC (partial rank correlation coefficient)} is the dominant
global SA method in mathematical epidemiology.
It is computationally cheap, handles monotone (but not necessarily linear)
relationships, and the Marino et al.~\cite{marino2008methodology} MATLAB implementation
is widely used and freely available; it builds on the rank-correlation
approach of Iman and Conover~\cite{iman1982distribution} and the
epidemiological SA tradition of Blower and Dowlatabadi~\cite{blower1994sensitivity}.

\subsection{The PRCC Calculation}

Given $N$ LHS samples of the $m$ \POIs and the corresponding $N$
model outputs:

\textbf{Step 1.} Rank-transform all inputs and the output separately.
Replace each value with its rank among the $N$ observations.

\textbf{Step 2.} Fit a linear regression of the rank-transformed output
on all rank-transformed inputs.

\textbf{Step 3.} For each parameter $p_j$, compute the partial correlation
coefficient between the rank-transformed $p_j$ and the rank-transformed output,
controlling for all other parameters. This is the PRCC for $p_j$.

\textbf{Step 4.} Compute a $t$-test $p$-value for each PRCC to assess
statistical significance.

The PRCC ranges from $-1$ to $+1$. Its sign has the same interpretation
as the SI sign rule from Chapter~2: positive means POI and QOI increase
together; negative means they move in opposite directions.

\subsection{When PRCC is Appropriate}

PRCC requires monotone relationships: the QOI must consistently increase
(or consistently decrease) as each POI increases, across the full
range of parameter values. It does not require linearity.
For most compartmental epidemic models, the relationships are approximately
monotone, which is why PRCC dominates the field.

Before applying PRCC, verify the monotonicity assumption by producing a
scatter matrix: plot the model output against each \POI for the full
LHS sample using MATLAB's \texttt{plotmatrix} function.
A roughly monotone cloud (upward or downward sloping, without reversals)
supports the use of PRCC. A hump-shaped or U-shaped cloud signals
non-monotonicity, and Sobol indices should be used instead.

PRCC and Sobol indices are not equivalent. PRCC assumes monotone
relationships; Sobol indices make no such assumption. When a relationship
is nonmonotone (a QOI first increases then decreases as a parameter
increases), PRCC may give a near-zero coefficient even for a highly
important parameter, while the Sobol total index will correctly identify
it as important. For the models in this course, both methods typically
agree, but verifying agreement is a useful consistency check.

%%------------------------------------------------------------
\section{Sobol Decomposition and Variance-Based SA}
\label{sec:sobol}
%%------------------------------------------------------------

Variance-based SA is the most rigorous framework for global SA:
it quantifies exactly how much of the total output variance is
attributable to each \POI or group of \POIs~\cite{sobol2001global}.

\subsection{The Intuition: Conditional Expectation}

Suppose parameter $Q_1$ is fixed at $q_1$ while all other parameters
vary across their ranges. The expected output $\E[Y\,|\,Q_1 = q_1]$
is some function of $q_1$. If $Q_1$ is important, this conditional
expectation varies substantially as $q_1$ changes.
The variance of this conditional expectation,
$\Var_{Q_1}[\E_{Q_{-1}}(Y\,|\,Q_1)]$, measures how much of $Y$'s
variability is explained by $Q_1$ alone. Normalizing by the total
variance gives the first-order Sobol index.

\subsection{The ANOVA Decomposition}\index{ANOVA decomposition}

\textbf{Why does a unique decomposition exist?}
Because statistical independence allows us to cleanly separate
what each parameter contributes alone from what it contributes jointly.
When the $Q_i$ are independent, projecting $f$ onto the subspace
of functions of $Q_i$ alone gives a well-defined component $f_i(Q_i)$,
and these projections do not interfere with each other.
The orthogonality $\int_0^1 f_i(Q_i)\,dQ_i = 0$ is the mathematical
expression of this separation: each component has zero mean after
integrating out $Q_i$, so distinct components contribute zero covariance.
When parameters are correlated, this orthogonality breaks down,
and the decomposition loses its uniqueness
(Section~\ref{sec:correlatedsobol} returns to this point).
With independence, the decomposition is both unique and additive in variance.

For independent \POIs $\vec{Q} = (Q_1,\ldots,Q_m)$ uniformly
distributed on $[0,1]^m$, any square-integrable model $Y = f(\vec{Q})$
admits the unique \textbf{ANOVA decomposition}:
\begin{equation}
  f(\vec{Q}) = f_0 + \sum_i f_i(Q_i)
             + \sum_{i<j} f_{ij}(Q_i, Q_j) + \cdots + f_{12\cdots m}(\vec{Q}),
  \label{eq:ANOVA}
\end{equation}
where $f_0 = \E[Y]$, $f_i(Q_i) = \E[Y\,|\,Q_i] - f_0$, and the
higher-order terms capture interactions. Uniqueness follows from the
orthogonality conditions $\int_0^1 f_{u}(\vec{Q}_u)\,dQ_i = 0$
for each $i \in u$; see Sobol~\cite{sobol2001global} Theorem~1
and Saltelli et al.~\cite{saltelli2008global} Appendix~A for a proof.
The terms are orthogonal, so their variances add:
\begin{equation}
  D = \Var(Y) = \sum_i D_i + \sum_{i<j} D_{ij} + \cdots + D_{12\cdots m},
  \label{eq:variancedec}
\end{equation}
where $D_i = \Var[\E(Y\,|\,Q_i)]$ and $D_{ij} = \Var[\E(Y\,|\,Q_i,Q_j)] - D_i - D_j$.

\medskip\noindent
\textbf{Worked example: a model with an interaction term.}
Consider $f(Q_1, Q_2) = Q_1 + Q_2 + Q_1 Q_2$ with $Q_1, Q_2 \sim \mathcal{U}(0,1)$
independent. Compute the ANOVA components directly from the definitions:
\begin{align*}
  f_0 &= \E[Y] = \tfrac{1}{2} + \tfrac{1}{2} + \tfrac{1}{4} = \tfrac{5}{4},\\
  f_1(Q_1) &= \E[Y\,|\,Q_1] - f_0
    = \bigl(Q_1 + \tfrac{1}{2} + \tfrac{Q_1}{2}\bigr) - \tfrac{5}{4}
    = \tfrac{3Q_1}{2} - \tfrac{3}{4} = \tfrac{3}{2}\bigl(Q_1 - \tfrac{1}{2}\bigr),\\
  f_2(Q_2) &= \E[Y\,|\,Q_2] - f_0
    = \tfrac{3}{2}\bigl(Q_2 - \tfrac{1}{2}\bigr) \quad\text{(by symmetry)},\\
  f_{12}(Q_1,Q_2) &= f - f_0 - f_1 - f_2
    = Q_1 Q_2 - \tfrac{Q_1}{2} - \tfrac{Q_2}{2} + \tfrac{1}{4}
    = \bigl(Q_1 - \tfrac{1}{2}\bigr)\bigl(Q_2 - \tfrac{1}{2}\bigr).
\end{align*}
Verify orthogonality: $\int_0^1 f_{12}\,dQ_1 = \bigl(\int_0^1(Q_1-\tfrac{1}{2})\,dQ_1\bigr)(Q_2-\tfrac{1}{2}) = 0$. \checkmark

The variances are $D_1 = D_2 = \tfrac{3}{4} \cdot \Var(Q_1) = \tfrac{3}{4}\cdot\tfrac{1}{12} = \tfrac{1}{16}$
and $D_{12} = \Var(f_{12}) = \Var(Q_1-\tfrac{1}{2})\Var(Q_2-\tfrac{1}{2}) = \tfrac{1}{144}$.
Total variance $D = \tfrac{1}{16} + \tfrac{1}{16} + \tfrac{1}{144} = \tfrac{19}{144}$.
The first-order indices $S_1 = S_2 = \tfrac{1/16}{19/144} = \tfrac{9}{19} \approx 0.47$,
and the interaction index $S_{12} = \tfrac{1/144}{19/144} = \tfrac{1}{19} \approx 0.053$.
Note $S_1 + S_2 + S_{12} = 1$. For the additive model $f = Q_1 + Q_2$
(no interaction), $S_1 = S_2 = \tfrac{1}{2}$ and $S_{12} = 0$: the indices
sum to 1 with nothing left over for interactions.
This is why $\sum S_i < 1$ signals the presence of interactions.

\subsection{Sobol Indices}

\begin{definition}[Sobol Indices]\index{Sobol indices}
The \textbf{first-order Sobol index}\index{Sobol indices!first-order} for parameter $i$ is:
\begin{equation}
  S_i = \frac{D_i}{D} = \frac{\Var[\E(Y\,|\,Q_i)]}{\Var(Y)}.
  \label{eq:S1}
\end{equation}
The \textbf{total Sobol index}\index{Sobol indices!total} for parameter $i$~\cite{homma1996importance} is:
\begin{equation}
  S_{T_i} = 1 - \frac{\Var[\E(Y\,|\,\vec{Q}_{-i})]}{\Var(Y)}
           = \frac{\E[\Var(Y\,|\,\vec{Q}_{-i})]}{\Var(Y)},
  \label{eq:ST}
\end{equation}
where $\vec{Q}_{-i}$ denotes all parameters except $Q_i$.
\end{definition}

The first-order index measures the \emph{direct} effect of $Q_i$;
the total index measures the direct effect plus all interactions with
other parameters. When $S_{T_i} = S_i$ for all $i$, the model is
additively separable (no interactions). When $S_{T_i} > S_i$,
parameter $i$ interacts with other parameters.

\unifyingidea{1} The first-order Sobol index $S_i = D_i/D$ uses the
same ratio structure as the sensitivity index from Chapter~2;
it measures one variance as a fraction of another. The connection
is explicit: for a linear model $Y = \sum_i a_i Q_i$ with independent
$Q_i \sim \mathcal{U}[l_i, u_i]$, the first-order Sobol index equals
the square of the sigma-normalized index: $S_i = (S_i^\sigma)^2$.
Global SA generalizes local SA; for linear models, they agree exactly.

\unifyingidea{3} The Sobol decomposition resolves the tension that
Unifying Idea~3 created in Chapter~8. Local SA is a linearization,
valid near the nominal values. Sobol indices make no linearity
assumption: they measure the actual variance contribution across
the full parameter range. When the linearization is poor (Model~B
from the chapter opening), local and global SA give different answers,
and Sobol indices give the correct answer.

For readers with a functional analysis background (Appendix~\ref{app:innerproduct}), the ANOVA decomposition~\eqref{eq:ANOVA}
is the orthogonal decomposition of $f$ in the Hilbert space
$L^2([0,1]^m)$ with respect to the inner product
$\langle f, g\rangle = \int_{[0,1]^m} f(\vec{q})\,g(\vec{q})\,d\vec{q}$.
The orthogonality explains both why the variances add
in~\eqref{eq:variancedec} and why the method requires independent \POIs:
independence is needed for the inner product to factorize and the
projections to be orthogonal.

\begin{selfcheck}
For $Y = c_1 Q_1 + c_2 Q_2$ with $Q_i \sim \mathcal{U}(0,1)$ independent:
(a) Compute $\E[Y\,|\,Q_1]$ and $\Var[\E(Y\,|\,Q_1)]$.
(b) Compute $\Var(Y)$ using the independence of $Q_1$ and $Q_2$.
(c) Show that $S_1 = c_1^2/(c_1^2 + c_2^2)$ and $S_2 = c_2^2/(c_1^2 + c_2^2)$.
(d) Verify $S_1 + S_2 = 1$ (no interaction for additive models).
\end{selfcheck}
\begin{scAnswer}
(a) $\E[Y\,|\,Q_1] = c_1 Q_1 + c_2/2$.
$\Var[\E(Y\,|\,Q_1)] = c_1^2 \Var(Q_1) = c_1^2/12$.

(b) $\Var(Y) = c_1^2\Var(Q_1) + c_2^2\Var(Q_2) = (c_1^2 + c_2^2)/12$.

(c) $S_1 = D_1/D = (c_1^2/12)/((c_1^2+c_2^2)/12) = c_1^2/(c_1^2+c_2^2)$. Similarly $S_2$.

(d) $S_1 + S_2 = (c_1^2 + c_2^2)/(c_1^2+c_2^2) = 1$. For linear additive models,
there are no interaction terms and first-order indices sum to 1.
\end{scAnswer}

%%------------------------------------------------------------
\section{Computing Sobol Indices: The Saltelli Algorithm}
\label{sec:saltellialg}
%%------------------------------------------------------------

A naive Monte Carlo estimate of $S_i$ requires choosing $Q_i = q_i$
for many values and averaging the model over all other parameters
for each, this costs $O(N^2)$ evaluations. The insight of Saltelli~\cite{saltelli2002making}
is that all first-order and total Sobol indices can be estimated with
$N(m+2)$ evaluations using two carefully constructed sample matrices.
The design and estimators for the total sensitivity index are developed
in Saltelli et al.~\cite{saltelli2010variance}.

\subsection{The Saltelli Construction}

Generate two independent LHS matrices $A$ and $B$, each $N\times m$.
Define the $m$ matrices $A_B^{(i)}$: matrix $A$ with its $i$-th column
replaced by the $i$-th column of $B$. The model is evaluated at all
rows of $A$, $B$, and all $m$ matrices $A_B^{(i)}$: a total of
$N(m+2)$ evaluations.

\subsection{The Jansen Estimators}\index{Jansen estimators}

For the first-order index, use the Jansen~\cite{jansen1999analysis} estimator:
\begin{equation}
  \hat{S}_i = 1 - \frac{\frac{1}{2N}\sum_{j=1}^N [f(B_j) - f(A_{B,j}^{(i)})]^2}{\hat{D}},
  \label{eq:Jansen_S1}
\end{equation}
where $\hat{D} = \frac{1}{N}\sum_j f(A_j)^2 - \left(\frac{1}{N}\sum_j f(A_j)\right)^2$.

For the total index:
\begin{equation}
  \hat{S}_{T_i} = \frac{\frac{1}{2N}\sum_{j=1}^N [f(A_j) - f(A_{B,j}^{(i)})]^2}{\hat{D}}.
  \label{eq:Jansen_ST}
\end{equation}
Use these Jansen estimators, not the original Saltelli (2002) formulas,
which have higher variance. The Jansen estimator is consistent
(converges to the true $S_i$ as $N\to\infty$) and nearly unbiased;
see Jansen~\cite{jansen1999analysis} Theorem~1.
The key step is that
\[
  \E\!\left[\frac{1}{2N}\sum_{j=1}^N (f(B_j) - f(A_{B,j}^{(i)}))^2\right]
  = D - D_i = \Var(Y) - \Var[\E(Y\,|\,Q_i)],
\]
which follows because $B_j$ and $A_{B,j}^{(i)}$ share all parameter
values except column $i$ (from $A$), so their difference isolates
the $Q_i$-dependent variation. Dividing by $\hat{D}$ and rearranging
gives $1 - \hat{S}_i \approx (D - D_i)/D = 1 - S_i$.
The estimator is nearly unbiased at finite $N$ because the bias
is $O(1/N)$~\cite{jansen1999analysis}, negligible for $N \geq 1000$.

\subsection{Convergence}

The estimators converge as $O(1/\sqrt{N})$, the standard Monte Carlo rate.
In practice, $N \geq 1000$ is needed for reliable estimates.
For $m = 4$ parameters (SIR model): $N(m+2) = 6000$ evaluations at $N = 1000$.
If each evaluation takes 0.1~seconds: 10~minutes.

Plotting $\hat{S}_i$ and $\hat{S}_{T_i}$ vs.\ $N$ reveals the convergence:
estimates are noisy for $N < 500$ and stable for $N \geq 2000$.
This convergence plot is the diagnostic that answers
the misconception ``more samples always help'': below a threshold
$N$, additional samples do not produce reliable estimates because
the variance of the estimator exceeds the signal being measured.

\textbf{Confidence intervals via bootstrap.}
A point estimate $\hat{S}_i$ is useful but incomplete: it gives no
information about sampling uncertainty.
The standard approach is the \textbf{bootstrap}: resample the
$N(m+2)$ model evaluations with replacement, recompute $\hat{S}_i$
for each resample, and take the 2.5th and 97.5th percentiles of
the resulting distribution as a 95\% confidence interval.
In MATLAB:
\begin{verbatim}
B = 500;          %% number of bootstrap replicates
S1_boot = zeros(B, m);
idx_all = 1:N;
for b = 1:B
    idx = idx_all(randi(N, N, 1));     %% resample rows
    A_b = A(idx,:);  B_b = Bmat(idx,:);
    [S1_boot(b,:), ~] = sobol_jansen(@(x) model_eval(x), A_b, B_b);
end
S1_CI = prctile(S1_boot, [2.5 97.5]);   %% 95% CI for each index
\end{verbatim}
A confidence interval that straddles zero indicates that the parameter's
effect cannot be distinguished from noise at the current sample size.
A confidence interval that excludes zero confirms the parameter is
genuinely influential. As a practical rule: if the confidence interval
width exceeds the index value, double $N$ before drawing conclusions.
When $m$ is large, however, even a single Sobol run at $N = 2000$
may be prohibitively expensive; Morris screening addresses exactly this case.

%%------------------------------------------------------------
\section{Morris Screening}\index{Morris screening}
\label{sec:morris}
%%------------------------------------------------------------

When $m$ is large (30 or more parameters), running Sobol at $N = 2000$
costs $N(m+2) = 64{,}000$ model evaluations, potentially prohibitive.
Morris screening~\cite{morris1991factorial} provides a cheap first step:
$N(m+1)$ evaluations (typically $N = 10$ trajectories) to identify
which parameters are truly unimportant. An improved radial design
that gives more uniform coverage with fewer trajectories is
described by Campolongo et al.~\cite{campolongo2007effective}.

For each trajectory, one parameter is changed at a time by a fixed
fraction $\Delta$ of its range. The \textbf{elementary effect}\index{elementary effect} of
parameter $i$ at point $\vec{q}$ is:
\begin{equation}
  EE_i(\vec{q}) = \frac{f(\vec{q} + \Delta\vec{e}_i) - f(\vec{q})}{\Delta}.
  \label{eq:EE}
\end{equation}
After $r$ trajectories, the mean $\mu_i^* = (1/r)\sum|EE_i|$
and standard deviation $\sigma_i$ of the elementary effects are computed.
Parameters with small $\mu_i^*$ are unimportant. Parameters with large
$\sigma_i$ have interactions or nonlinear effects.

The Morris diagnostic plot ($\sigma_i$ vs.\ $\mu_i^*$) separates
parameters into three groups: unimportant (small $\mu_i^*$, small $\sigma_i$);
important and linear (large $\mu_i^*$, small $\sigma_i$); important
and nonlinear or interacting (large $\mu_i^*$, large $\sigma_i$).

\textbf{Workflow}: run Morris screening first; fix unimportant parameters
at their nominal values; run Sobol only on the important subset.
This can reduce the Sobol cost by a factor of 5--10.

\medskip

\noindent\textbf{When even reduced Sobol is too expensive: surrogate models.}
If the model is computationally expensive and Morris screening reduces
the parameter set to, say, five influential \POIs, running Sobol at
$N = 2000$ still requires $N(m+2) = 14{,}000$ model evaluations.
For a climate model where each run takes hours, this remains prohibitive.

The practical solution is a \textbf{surrogate model}\index{surrogate model}
(also called an emulator or meta-model): a cheap statistical
approximation to the expensive model, trained on a moderate number of
carefully chosen model runs.
Concretely, a Gaussian process surrogate fits a probability distribution
over functions to the training runs, predicting the model output at any
new input along with a confidence interval; a polynomial chaos expansion
instead represents the output as a truncated series in orthogonal polynomials
of the uncertain inputs, with coefficients estimated from the training set.
Gaussian process emulators and polynomial chaos expansions are the
most widely used surrogates. Once trained on roughly $50$--$200$
model evaluations, a surrogate can be evaluated millions of times
in seconds, making full Sobol estimation feasible even for
computationally expensive simulations.

Surrogate-based global SA does not replace the methods in this chapter;
it wraps them. The LHS sampling, the Saltelli construction, and the
Jansen estimators all apply to the surrogate exactly as they apply
to the original model. The surrogate simply makes the evaluations
fast enough to converge.
Smith~\cite{smith2014uncertainty} develops Gaussian process and
polynomial chaos surrogates in Chapters~12--13, building directly
on the LHS samples and Sobol indices from this chapter.
Students proceeding to the second semester will construct surrogates
using exactly the training sets generated in this chapter's laboratories.

%%------------------------------------------------------------
\section{Sobol Indices for Correlated POIs}\index{Sobol indices!correlated POIs}
\label{sec:correlatedsobol}
%%------------------------------------------------------------

Standard Sobol decomposition assumes independent \POIs. When parameters
are correlated, three problems arise: first-order indices can exceed 1 or
be negative (a diagnostic flag that correlation is present); total indices
may not sum to 1; and the ANOVA decomposition~\eqref{eq:ANOVA} is no longer
unique, because correlation destroys the orthogonality that guarantees uniqueness.

Three approaches address correlated \POIs in global SA:

\textbf{Grouped indices}: treat a cluster of correlated parameters as
a single group. Compute the group Sobol index (the fraction of variance
explained by the entire group jointly). This avoids allocating variance
between correlated parameters, which is not well-defined.

\textbf{Modified estimators}: Kucherenko et al.~\cite{kucherenko2012estimation}
extend the Saltelli algorithm to correlated inputs using copula-based
sampling. This requires knowledge of the joint distribution of the
correlated parameters.

\textbf{PRCC}: handles correlation naturally through the ``partial''
in partial rank correlation. When parameters are correlated, PRCC
is the most robust and accessible global SA method.

\textbf{Practical recommendation}: for models with correlated \POIs,
use PRCC as the primary global SA tool and grouped Sobol indices as
a secondary check. Report the correlation structure explicitly alongside
the SA results.

%%------------------------------------------------------------
\section{Comparison of All SA Methods}
\label{sec:methodcomparison}
%%------------------------------------------------------------

Table~\ref{tab:allSA} consolidates all SA methods from the book.
It supersedes the partial decision tables in earlier chapters.
For a broader treatment of global SA methods, see Saltelli et al.~\cite{saltelli2008global,saltelli2000sensitivity}
and the recent epidemiological applications in Lu and Borgonovo~\cite{lu2023global}.

\begin{table}[htb]
\centering
\renewcommand{\arraystretch}{1.25}
\begin{tabular}{p{2.4cm}p{1.3cm}p{3.5cm}p{2.2cm}p{2.3cm}}
\toprule
Method & Chapter & When to use & Assumption & Cost \\
\midrule
Local SI (analytic) & 2, 4 & Differentiable model; small uncertainty & Linearization valid & Cheap \\
Local SI (FD) & 3 & Black-box; small uncertainty & Linearization valid & $2mn_q$ evals \\
Adjoint & 5--6 & Differentiable; many \POIs; few \QOIs & Linearization valid & 2 solves \\
AD (reverse) & 7 & Differentiable code; $m \gg r$ & Linearization valid & 2 passes \\
PRCC & 9 & Monotone model; moderate uncertainty; corr.\ \POIs OK & Monotone & $N$ evals \\
Morris & 9 & Screening; $m \geq 30$ & None & $N(m+1)$ evals \\
Sobol & 9 & Nonlinear; large uncertainty; independent \POIs & Independent & $N(m+2)$ evals \\
Grouped Sobol & 9 & Correlated \POIs; known groups & Known groups & $N(m+2)$ evals \\
\bottomrule
\end{tabular}
\caption{Complete SA method comparison across all chapters.
         The choice of method depends primarily on whether the model
         is differentiable, whether uncertainty is small or large,
         and whether \POIs are independent. No single method is universally best.}
\label{tab:allSA}
\end{table}

\unifyingidea{2} The adjoint row in Table~\ref{tab:allSA} is the
payoff for Chapters~5 and 6. The adjoint is the transpose in whatever
space you are working in: matrix transpose for linear systems
(Chapter~5), reversed ODE for dynamical systems (Chapter~6), and
reversed chain rule for computation graphs (Chapter~7). Every
``2 solves'' entry in the adjoint row is the same unifying idea,
applied in a different mathematical space.

%%------------------------------------------------------------
\section{Bridge to Semester 2}
\label{sec:bridgesemester2}
%%------------------------------------------------------------

This book quantifies how model output uncertainty arises from
parameter uncertainty, assuming the model structure is correct.
The second semester, covering Smith's \emph{Uncertainty Quantification}
\cite{smith2014uncertainty}, asks harder questions: what if the model is wrong?
How do we fit models to data? How do we construct efficient surrogate
models to replace expensive simulations?
Smith's Chapters~7--8 develop Bayesian parameter estimation and Markov chain
Monte Carlo methods for fitting models to data, the natural sequel to the
identifiability analysis introduced in Chapter~3 of this book.

The connections between the two semesters are direct:

Smith Ch.~14 develops local SA at the graduate level, extending
Chapters~2--7 of this book with functional-analytic rigor.

Smith Ch.~15 develops global SA and Sobol indices at the graduate level,
extending Chapter~9 of this book with measure-theoretic foundations.

Smith Ch.~6 treats parameter selection using the sensitivity matrix;
the core tool introduced in Chapter~4 of this book.

Smith Chs.~12--13 develop model discrepancy and Gaussian process
surrogates, methods that reduce the cost of global SA by replacing
the expensive model with a cheap approximation. Surrogates are
built from the same LHS samples used to estimate Sobol indices
in this chapter.

Students who have completed this book have the foundational vocabulary,
the computational intuition, and the mathematical framework to engage
productively with the full UQ literature. The path from the sensitivity
index in Chapter~2 to the full Bayesian UQ framework in Smith is
a single, continuous intellectual development. As an example of
these methods applied to a realistic multi-parameter epidemic model,
see Qu, Xue, and Hyman~\cite{qu2018modeling}, which uses SA to
guide Wolbachia-based mosquito-borne disease control.

The student who has reached this point knows how to ask which parameters
matter, how to compute that answer efficiently and honestly, and, just
as importantly, when not to trust the answer.
Every mathematical model is an argument; sensitivity analysis is the
discipline that reveals which assumptions the argument depends on.

\medskip\noindent
\textit{Good sensitivity analysis does not end the discussion.
It tells you where to begin it.}

%%------------------------------------------------------------
\section{MATLAB Laboratory: Complete Global SA Pipeline for SIR}
\label{sec:ch9matlab}
%%------------------------------------------------------------

This laboratory implements the full global SA pipeline for the SIR model
and compares all four methods: sigma-normalized, PRCC, Sobol, and Morris.

\textbf{Setup.}
SIR model with four \POIs $(k, \beta, \tau, L)$ and response $J = I_{\max}$
(peak infectious count). Use the literature-based uncertainty ranges
from Chitnis et al.~\cite{chitnis2008determining}:
$k \in [2, 8]$, $\beta \in [0.02, 0.12]$, $\tau \in [4, 14]$ days,
$L \in [1000, 100000]$ days.

\textbf{Part 1: LHS and sigma-normalized indices.}
Draw $N = 500$ LHS samples using the \texttt{lhs\_sample} function.
Compute $\sigma_{J}$ from the samples. Compute the sigma-normalized index
for each \POI using the local SI from Chapter~2 and the sample
standard deviations. Report the ranking.

\textbf{Part 2: PRCC.}
Run the SIR model at all 500 LHS samples. Rank-transform all inputs
and the output. Compute PRCC using the \texttt{partialcorr} function
in MATLAB's Statistics and Machine Learning Toolbox
(or download Marino et al.'s implementation from the supplementary
material of \cite{marino2008methodology}).
Produce a PRCC tornado plot and report $p$-values.

\textbf{Part 3: Sobol estimation.}
Draw $N = 2000$ samples using the quasi-random Halton or Sobol sequence
(MATLAB: \texttt{haltonset} or \texttt{sobolset}).
Build the $A$, $B$, and $m$ matrices $A_B^{(i)}$: $2000\times6 = 12{,}000$
total evaluations.
Compute $\hat{S}_i$ and $\hat{S}_{T_i}$ using the Jansen estimators.

\begin{verbatim}
%% From: src/globalsa/sobol_jansen.m (see GitHub repository)
rng(42);   %% reproducibility
[A, B]      = saltelli_sample(N, m, lb, ub);
[S1, ST, CI_S1] = sobol_jansen(@model_eval, A, B, 500, 0.05);
%% CI_S1(:,j) = [lower; upper] 95% confidence interval for S1(j)
\end{verbatim}

\textbf{Part 4: Morris screening.}
Implement $r = 10$ Morris trajectories with $\Delta = 0.4$.
Compute $\mu_i^*$ and $\sigma_i$ for each parameter.
Produce the Morris diagnostic plot ($\sigma_i$ vs.\ $\mu_i^*$).
Which parameter(s) could be fixed without significant loss of accuracy?

\textbf{Part 5: Convergence check.}
Plot $\hat{S}_\beta$ and $\hat{S}_{T_\beta}$ vs.\ $N$ for
$N \in \{100, 200, 500, 1000, 2000, 5000\}$.
At what $N$ do the estimates stabilize to within $\pm0.05$?

\textbf{Part 6: Full comparison table.}
Fill in the following table from your results:

\begin{center}
\begin{tabular}{lcccc}
\toprule
Parameter & Local SI & $S_i^\sigma$ & PRCC & $S_i$ (Sobol) \\
\midrule
$k$ & & & & \\
$\beta$ & & & & \\
$\tau$ & & & & \\
$L$ & & & & \\
\bottomrule
\end{tabular}
\end{center}

Where do the four methods agree? Where do they disagree? What accounts
for the differences?

%%------------------------------------------------------------
\section{Exercises}
%%------------------------------------------------------------

\begin{exercise}[\bloom{L3}]
\label{ex:ch9:sigmanorm}
For the SEIR model with local SIs $\SI{\tau_E}{R_0} = 0.6$,
$\SI{\tau_I}{R_0} = 1.2$, $\SI{\beta}{R_0} = 1.0$, $\SI{k}{R_0} = 1.0$
and uncertainty ranges $\sigma_{\tau_E} = 2$~days, $\sigma_{\tau_I} = 3$~days,
$\sigma_\beta = 0.01$, $\sigma_k = 0.8$:
\begin{enumerate}[label=(\alph*)]
  \item Compute the output standard deviation $\sigma_{R_0}$ using the
        delta method, at nominal $R_0^* = 2.0$.
  \item Compute all four sigma-normalized indices.
  \item Produce a tornado plot of the sigma-normalized indices.
  \item Compare the ranking with the local SI ranking. Which parameter
        changed rank most dramatically? Why?
\end{enumerate}
\end{exercise}

\begin{exercise}[\bloom{L3}]
\label{ex:ch9:LHS}
\begin{enumerate}[label=(\alph*)]
  \item Implement the \texttt{lhs\_sample} function. For $N = 100$,
        $m = 3$, $[0,1]^3$, generate an LHS and a random sample of the
        same size.
  \item Produce scatter plots of all three pairs of parameters for
        both samples (6 plots total).
  \item Compute the mean and standard deviation of the sample along each
        dimension. Which method is closer to the true uniform
        mean and standard deviation?
  \item For one parameter, plot the histogram of LHS vs.\ random
        sampling. Which covers the range more evenly?
\end{enumerate}
\end{exercise}

\begin{exercise}[\bloom{L3}]
\label{ex:ch9:PRCC}
For the SIR model with parameter ranges from the laboratory setup:
\begin{enumerate}[label=(\alph*)]
  \item Generate $N = 200$ LHS samples and run the SIR model at each.
  \item Compute PRCC for all four parameters with respect to $I_{\max}$.
  \item Produce a PRCC bar chart sorted by $|$PRCC$|$.
  \item Which parameter has the strongest negative PRCC? What does this
        mean in plain language?
  \item Compare the PRCC ranking with the local SI tornado plot from
        Chapter~2. Do they agree?
\end{enumerate}
\end{exercise}

\begin{exercise}[\bloom{L4}]
\label{ex:ch9:sobolanalytic}
For $Y = 2Q_1^2 + Q_2$ with $Q_1, Q_2 \sim \mathcal{U}(0,1)$, independent:
\begin{enumerate}[label=(\alph*)]
  \item Compute $\E[Y]$, $\Var(Y)$, $D_1 = \Var[\E(Y\,|\,Q_1)]$,
        $D_2 = \Var[\E(Y\,|\,Q_2)]$.
  \item Compute $S_1$, $S_2$. Do they sum to 1? Explain.
  \item Compute the total indices $S_{T_1}$, $S_{T_2}$.
  \item Now add $Q_1$ and $Q_2$ to get $Y' = 2Q_1^2 + Q_2 + Q_1 Q_2$
        (an interaction term). Without computing: does adding the interaction
        term change $S_1$, $S_2$, or both? In which direction?
        Then verify by computing the Sobol indices for $Y'$ analytically.
\end{enumerate}
\end{exercise}

\begin{exercise}[\bloom{L4}]
\label{ex:ch9:sobolSIR}
Using the Jansen estimator implementation from the laboratory:
\begin{enumerate}[label=(\alph*)]
  \item Compute $\hat{S}_i$ and $\hat{S}_{T_i}$ for all four SIR parameters
        at $N = 2000$.
  \item For which parameters is $S_{T_i} \gg S_i$? What does this indicate
        about their role in the model?
  \item Compute $\sum_i S_i$ and $\sum_i S_{T_i}$. What do you expect
        these to sum to for an additive model? Are your results consistent
        with this expectation?
  \item If you could measure exactly one parameter very precisely
        (reducing its uncertainty to zero), which parameter would reduce
        $\Var(I_{\max})$ the most? Justify using the Sobol indices.
\end{enumerate}
\end{exercise}

\begin{exercise}[\bloom{L4}]
\label{ex:ch9:morris}
For the SIR model, implement $r = 15$ Morris trajectories with $\Delta = 1/3$:
\begin{enumerate}[label=(\alph*)]
  \item Compute $\mu_i^*$ and $\sigma_i$ for each parameter.
  \item Produce the Morris diagnostic plot.
  \item Based on the Morris results, which parameter appears most influential?
        Which appears least influential?
  \item Run a reduced Sobol analysis using only the three most important
        parameters identified by Morris (fixing $L$ at its nominal value).
        Compare the cost and results with the full four-parameter Sobol.
\end{enumerate}
\end{exercise}

\begin{exercise}[\bloom{L5}]
\label{ex:ch9:globaldesign}
Choose a published epidemic model from the literature with at least
8 parameters. Design a complete global SA study:
\begin{enumerate}[label=(\alph*)]
  \item State the model, its parameters, and a justification for the
        uncertainty ranges you will use (with citations).
  \item State the \QOI or \QOIs and explain why they are the most
        relevant for the model's purpose.
  \item Choose between PRCC and Sobol as the primary global SA method.
        Justify your choice based on the model's properties.
  \item Specify the sample size $N$ and justify it based on convergence.
  \item Describe how you would present the results to a non-technical
        audience (e.g., a public health committee or funding agency).
  \item Identify one situation where your global SA results could still
        be misleading, and describe what additional analysis you would
        recommend.
\end{enumerate}
\end{exercise}

\begin{exercise}[\bloom{L6}]
\label{ex:ch9:evaluation}
A published SA study reports Sobol first-order and total indices for
a 15-parameter water-quality model. The authors note in the data section
that three parameters, phosphorus loading rate, nitrogen loading rate,
and water residence time, are positively correlated across the sampled
watersheds. The SA section reports standard Sobol indices computed using
independent sampling for all 15 parameters, without addressing the
correlation.
\begin{enumerate}[label=(\alph*)]
  \item Identify the specific mathematical problem that correlation
        creates for standard Sobol indices.
  \item For each of the three correlated parameters, state whether the
        reported first-order index is likely over-estimated, under-estimated,
        or unchanged, and explain why.
  \item Propose the minimal methodological correction that would make
        the analysis valid.
  \item The authors justify their approach by noting that their Sobol
        indices sum to approximately 1.0. Explain why this is not
        a valid check for independence.
\end{enumerate}
\end{exercise}

%%------------------------------------------------------------

